\textbf{结论名称：} 面面平行判定定理

\textbf{适用条件：} 空间中有两个平面 $\alpha$、$\beta$，以及位于平面 $\alpha$ 内的两条相交直线 $a$、$b$，且这两条直线分别平行于平面 $\beta$。

\textbf{变量说明：} 
$a,b$：平面 $\alpha$ 内的两条直线； 
$\alpha,\beta$：空间中的两个平面； 
$P$：直线 $a$ 与 $b$ 的交点。

\textbf{核心公式：}
\[
\boxed{\left.\begin{array}{l}
a \subset \alpha,\; b \subset \alpha\\
a \cap b = P\\
a \parallel \beta,\; b \parallel \beta
\end{array}\right\}
\;\Longrightarrow\; \alpha \parallel \beta}
\]

\textbf{使用场景：}
用于证明两个平面平行。将“面面平行”转化为“线面平行”，再进一步转化为“线线平行”，在几何体截面平行证明、空间位置关系判定中频繁出现。

\textbf{简单例子：}
在长方体 $ABCD$-$A_1B_1C_1D_1$ 中，证明平面 $AB_1D_1 \parallel$ 平面 $C_1BD$。
由 $AB_1 \parallel C_1D$，且 $C_1D \subset$ 平面 $C_1BD$，得 $AB_1 \parallel$ 平面 $C_1BD$；同理 $AD_1 \parallel BC_1$，得 $AD_1 \parallel$ 平面 $C_1BD$。又 $AB_1$ 与 $AD_1$ 相交于点 $A$，且二者均位于平面 $AB_1D_1$ 内，根据面面平行判定定理可得平面 $AB_1D_1 \parallel$ 平面 $C_1BD$。